\chapter*{ЗАКЛЮЧЕНИЕ}
\addcontentsline{toc}{chapter}{ЗАКЛЮЧЕНИЕ}

По результатам выполненной курсовой работы сформулированы следующие выводы:

\begin{enumerate}
  \item В ходе проведённого аналитического обзора были рассмотрены основные методы количественного определения анионов в природных водах. После сравнения титриметрии, спектрофотометрии, ионной хроматографии и капиллярного электрофореза был сделан вывод о том, что капиллярный электрофорез является наиболее целесообразным методом для выполнения поставленной задачи за счёт низкого расхода реагентов, небольшого времени анализа благодаря возможности проведения многокомпонентного анализа и высокой эффективности разделения. Для количественного определения анионов без хромофорных групп использовалась модификация с косвенным УФ-детектированием.

  \item Была разработана методология проведения эксперимента, основанная на системе капиллярного электрофореза «Капель-105М». Также были установлены и подтверждены рабочие условия анализа: длина волны детектирования 254~нм, рабочее напряжение $-25$~кВ, температура термостатирования 20~$^\circ$C и состав фонового электролита — оксид хрома (VI) ($10$~ммоль/дм$^3$), диэтаноламин ($30$~ммоль/дм$^3$) и ЦТА-OH ($2$~ммоль/дм$^3$). Были построены линейные градуировочные зависимости для исследуемых анионов. Полученные пределы количественного и качественного обнаружения находились в диапазонах 0,03-0,15~мг/дм$^3$ и соответствовали требованиям нормативных документов по измерению качества природных вод.

  \item Был проведён анализ реальных поверхностных вод реки Прохладной в четырёх точках отбора проб. Для каждой пробы было выполнено по три параллельных измерения. Относительное стандартное отклонение не превышало 1,6\%, а доверительный интервал составлял 3,9\%. Содержание хлоридов, сульфатов, нитритов и нитратов во всех точках не превышало установленных ПДК, но несмотря на это зафиксировано превышение рыбохозяйственного ПДК по ортофосфатам в 2–3 раза во всех четырёх пробах, что свидетельствует об экологической фосфатной нагрузке на исследованный водный объект и повышает риск развития эвтрофикационных процессов.
\end{enumerate}

Перспективы дальнейшего мониторинга целесообразно реализовывать по нескольким направлениям. Первым направлением является расширение количества измеряемых компонентов природных вод, например, включение в анализ бромид-ионов и органических кислот, что позволит получить более полную картину экологического состояния природных вод. Вторым направлением развития стоит избрать проведение сезонных анализов воды на нескольких участках, что позволит оценить динамику фосфатной нагрузки и выявить её связь с гидрологическим режимом и хозяйственной деятельностью.
